%------------------------------------------------------------------------------%
\section{Equações e Inequações modulares}

%------------------------------------------------------------------------------%
\subsection{Soluções de Equações Modulares}

Dadas as expressões algébricas $A$ e $B$, as
soluções da equação $\abs{A}=B$ devem satisfazer
\[
  B \geq 0
  \quad\text{ e }\quad
  \left(
    A = B
    \quad\text{ ou }\quad
   -A = B
  \right).
\]

%------------------------------------------------------------------------------%
\begin{example}
  Resolva a equação modular
  \(\abs{5x-2} = 13\).

  \solution
  Observe que nesse caso temos uma equação do tipo $\abs{A}=B$, com $B>0$.

  Logo,
  \[
    5x-2 = 13
    \quad\text{ ou }\quad
    -(5x-2)=13,
  \]
  resolvendo essas duas equações obtemos que o conjunto solução da
  equação pedida é
  \[
    S = \left\{-\frac{11}{5},\, 3 \right\}.
  \]
\end{example}
%------------------------------------------------------------------------------%

%------------------------------------------------------------------------------%
\begin{example}
  Resolva a equação modular
  \(\abs{2x-5} - \abs{x-2} = 3x+1\).

  \solution
  Para resolver essa equação modular, dividiremos a análise dos
  módulos em etapas, analisando cada módulo separadamente.

  Primeiro caso $\abs{2x-5}$
  \[
    |2x-5|=  2x-5 , se 2x-5 \geq 0 \to x \geq \frac{5}{2}.
  \]
  \[
    |2x-5|=-(2x-5), se 2x-5 <    0 \to x    < \frac{5}{2}.
  \]

  Segundo caso $\abs{x-2}$
  \[
    |x-2|=  x-2 , se x-2 \geq 0 \to x \geq 2.
  \]
  \[
    |x-2|=-(x-2), se x-2 <    0 \to x <    2.
  \]

  Fazendo a análise nos intervalos
  \begin{ceqn}
  \[
    \begin{array}{c<{\;}>{\;}c<{\;}>{\;}c<{\;}>{\;}c} \hline
      & \displaystyle\left(-\infty, 2\right)
      & \displaystyle\left[2, \sfrac{5}{2}\right)
      & \displaystyle\left[\sfrac{5}{2}, \infty\right) \\ \hline
      \abs{2x-5}           & -(2x-5) & -(2x-5) & 2x-5 \\
      \abs{x-2}            & -(x-2)  &   x-2   &  x-2 \\
      \abs{2x-5}-\abs{x-2} &  -x+3   &  -3x+7  &  x-3 \\ \hline
    \end{array}
  \]
  \end{ceqn}

  Intervalo $(-\infty, 2)$
  \begin{align*}
    -x+3 & = 3x+1        \\
    -4x  & = -2          \\
    x    & = \frac{1}{2}
  \end{align*}
  Note que \(x= \sfrac{1}{2}\) pertence ao intervalo $(-\infty, 2)$.

  Intervalo \(\left[2, \sfrac{5}{2}\right)\)
  \begin{align*}
    -3x+7 & = 3x+1 \\
    -6x   & = -6   \\
    x     & = 1
  \end{align*}
  Note que $x = 1$ não pertence ao intervalo
  \(\left[2, \sfrac{5}{2}\right)\),
  logo não pode ser solução da equação.

  Intervalo \(\left[\sfrac{5}{2}, \infty\right)\)
  \begin{align*}
    x-3 & = 3x+1 \\
    -2x & = 4    \\
    x   & = -2
  \end{align*}
  Note que $x= -2$ não pertence ao intervalo
  \(\left[\sfrac{5}{2}, \infty\right)\),
  logo não pode ser solução da equação.

  Portanto a única solução da equação é
  \(
    S = \left\{ \sfrac{1}{2} \right\}
  \).
\end{example}

%------------------------------------------------------------------------------%
\subsection{Soluções de Inequações Modulares}

Dada a constante real $c \geq 0$ e a expressão algébrica $A$, temos
\begin{center}
  \vspace{-0.5\baselineskip}
  \begin{tabular}{c<{\;}>{\;}c}
    \hline
     Inequação        & Condição equivalente                          \\ \hline
    \(\abs{A}\leq c\) & \(-c \leq A \leq c\)                          \\
    \(\abs{A}\geq c\) & \( A \leq -c \quad\text{ ou }\quad A \geq c\) \\ \hline
  \end{tabular}
\end{center}

%------------------------------------------------------------------------------%
\begin{example}
  Resolva a inequação \(\abs{2x+3} \geq 4\).

  \solution
    Resolver a inequação
    \[
      \abs{2x+3} \geq 4
    \]
    é equivalente resolver as inequações
    \[
      2x+3 \leq -4
      \quad\text{ e }\quad
      2x+3 \geq 4.
    \]
    Resolvendo a primeira inequação
    \[
      2x+3 \leq -4 \Rightarrow 2x \leq -7 \Rightarrow x  \leq -\frac{7}{2}.
    \]
    Resolvendo a segunda inequação
    \[
      2x+3 \geq 4 \Rightarrow 2x \geq 1 \Rightarrow x \geq \frac{1}{2}.
    \]
    Fazendo a união dos intervalos encontrados obtemos que a solução da
    inequação pedida é
    \[
      \left\{ x \in \R \st
              x \leq -\frac{7}{2}
              \;\text{ ou }\;
              x \geq \frac{1}{2}
      \right\}.
    \]
\end{example}
%------------------------------------------------------------------------------%

%------------------------------------------------------------------------------%
\begin{example}
  Resolva a inequação \(\abs{4x^2-1} \leq x+2\).

  \solution
  Nesse caso, primeiro temos que impor que $x+2 \geq 0$, isto é $x \geq -2$.
  Resolver a inequação $\vert 4x^2-1 \vert \leq x+2$ é equivalente resolver a
  inequação simultânea $-(x+2) \leq 4x^2-1 \leq x+2$.
  \begin{itemize}
      \item[1.] $-x-2 \leq 4x^2-1 \Rightarrow 4x^2+x+1 \geq 0$, que é uma
        inequação do segundo grau (Veja no Cap.2 como resolver), cuja solução
        é todo o conjunto dos números reais, isto é, $S_1= \{ x \in
        \R \}$.

      \item[2.] $4x^2-1 \leq x+2 \Rightarrow 4x^2-x-3 \leq 0$, cuja
        solução é $S_2= \{ x \in \R| -\frac{3}{4} \leq x \leq 1 \}$.
  \end{itemize}
  Fazendo a intersecção desses intervalos, temos que a solução geral da
  inequação modular é $S=\{ x \in \R | -\frac{3}{4} \leq x \leq 1
  \}$.
\end{example}
%------------------------------------------------------------------------------%

%------------------------------------------------------------------------------%
\begin{example}
  Resolva a inequação \(\abs{x-3} - \abs{5x-2} \geq x-2\).

  \solution
  {\large\color{red} Falta solução.}
\end{example}
%------------------------------------------------------------------------------%

%------------------------------------------------------------------------------%
